% Pursuing Stacks (cotes 134-2 to 134-8): archive cote 134-4, site page 16 (PDF page 17), typescript page 104.
% typescript. Draft compared: ps3.tex:793-853.
% AI-transcribed by Claude Opus 5.5 (claude-opus-5-5, adaptive thinking, effort high), Message Batches API, batch ps1 (msgbatch_01Rv5hCaNTtUt4JHx1yqoHPL), 2026-11-12; image review recorded in records/CHECK_DECISIONS_134_3_4.tsv.
\scanpage{134-4/16}
\mspage{104}
It is clear then,  for $f,g\hand{\in}\mathrm{Hom}(X,Y)$,  that any $\underline{I}$-homotopy from f to g
gives rise to a $\underline{J}$-homotopy; hence if f and g are $\underline{I}$-homotopic they are
$\underline{J}$-homotopic,  and hence the same for homotopy in the wide sense.  We get
thus a \inserted[typed]{canonical} functor
\[ M_{\underline{I}} \hand{\to} M_{\underline{J}} \]
which is the identity on objects,  entering into a cascade of canonical
functors
\[ M \hand{\to} M_{\underline{I}} \hand{\longrightarrow} M_{\underline{J}} \hand{\to} H \ , \]
\xout{where $H$ \unsure{is} $= H_M = W^{-1}M$}
where  $H=H_M= W^{-1} M$  can be viewed as the common localization of $M,M_{\underline{I}},M_{\underline{J}}$
with respect to the notion of weak equivalences in these categories.  We
may view $M_{\underline{J}}$  as a closer approximation to H than $M_{\underline{I}}$ .  T\over[typed]{H}{h}ere  is of course
an evident transitivity relation for the functors corresponding to two
composable morphisms of intervals
\[ \underline{K} \hand{\to} \underline{J} \hand{\to} \underline{I} \quad . \]
\uline{R\over[typed]{E}{e}mark}.  \over[typed]{I}{I}n order to get that $\underset{\underline{I}}{\hand{\sim}}$ implies $\underset{\underline{J}}{\hand{\sim}}$ ,  it is sufficient to
make a much weaker assumption than existence of a morphism of intervals
$\underline{J}\hand{\to}\underline{I}$  -  namely it suffices to assume that the two sections $\hand{\delta}^{\underline{I}}_{i}$ : e $\hand{\to} \underline{I}$
are $\underline{J}$-homotopic.  M$\overset{\hand{\unsure{\smile}}}{\underset{\hand{\wedge}}{\mathrm o}}$re generally,  let $\hand{\sim}$ be an equivalence relation in
F\over{l}{l}(M),  compatible with compositions and with cartesian products (this is
the case indeed for $\underset{\underline{J}}{\hand{\sim}}$ ),  and let  $\underline{I} \over{\unsure{b}}{=} (I,\hand{\delta}_0,\hand{\delta}_1)$ any object I \inserted[typed]{of M} (not neces-%
sarily aspheric over e) endowed with two sections over e ,  such that
$\hand{\delta}_0\hand{\sim}\hand{\delta}_1$ .   Then the interval $\underline{I}$ gives rise to an equivalence relation $\underset{\underline{I}}{\hand{\sim}}$
in F\over{l}{l}(M) , whose definition is quite independant of W -  and a priori
\xout{(if it were not for the condition} if  $f\underset{\underline{I}}{\hand{\sim}}g$  and $f\hand{\in}W$,  this need not
imply $g\hand{\in}W$.  However,  the condition $\hand{\delta}_0\hand{\sim}\hand{\delta}_1$  implies immediately that the
\xout{condition} relation  $\underset{\underline{I}}{\hand{\sim}}$  implies the relation $\hand{\sim}$ .  When the latter is
$\underset{\underline{J}}{\hand{\sim}}$ ,  we get moreover that W is the inverse image of a set of arrows in
$M_{\underline{I}}$ , \hand{i.e. $f\underset{I}{\sim}g$ and $f\in W$ implies $g\in W$.}

   An interesting particular case is the one when we can find a homotopy
interval  $\underline{I}_0$  in \over{M}{M} ,  which has the property that for any other homotopy
interval $\underline{I}$ in M,  it's structural sections satisfy
\[ \hand{\delta}^{\underline{I}}_0 \underset{\underline{I}_0}{\hand{\sim}} \hand{\delta}^{\underline{I}}_1 \qquad . \]
This implies that the homotopy relation $\underset{\underline{I}_0}{\hand{\sim}}$ is implied by all other simi-%
lar relations $\underset{\underline{I}}{\hand{\sim}}$ ,  i.e. it is the coarsest among all relations $\underset{\underline{I}}{\hand{\sim}}$ .  We
may \xout{call} then view $\underline{I}_0$ as a ``\uline{fundamental}'' or ``\uline{characteristic}'' \uline{homotopy
interval} in M,  in the sense that \xout{to check if} \inserted[typed]{the relation} $f\underset{W}{\hand{\sim}}g$  in the sense below ,
namely \deleted{\xout{if there exis\illegible{}}} \inserted[typed]{existence of} a homotopy  interval $\underline{I}$ such that $f\underset{\underline{I}}{\hand{\sim}}g$ ,  is equi-%
valent to  $f\underset{\underline{I}_0}{\hand{\sim}}g$  i.e. can be checked using \inserted[typed]{\over[typed]{T}{t}he} one and unique  $\underline{I}_0$ .  In the
case of M=(Cat\hand{)},  we get readily $\hand{\Delta}_1$ as a characteristic homotopy inter-
\margin{16}
